% GENERATED FILE. Edit canonical lessons, then run npm run generate:books.
% canonical-source-hash: fnv1a64:eb898972e449c56c
% canonical-lessons: UR-C237-mistri, UR-C237-bera, UR-C237-chaukidar, UR-C237-sipahi, UR-C237-fauj

\chapter{Mechanic, Waiter, Watchman, Soldier, Army}
\label{ch:ur-mechanic-waiter-watchman-soldier-army}
\hlchaptermodality{237}

\begin{chapteropening}
\textbf{By the end of this chapter:} \emph{I can say a mechanic, a waiter, a watchman, a soldier and an army.}

Complete the last lesson of chapter 237: I can say a mechanic, a waiter, a watchman, a soldier and an army.
\end{chapteropening}

\section[\texorpdfstring{\ur{مستری} (mistrī)}{مستری (mistrī)}]{\ur{مستری} (mistrī) — a mechanic}

\label{lesson:UR-C237-mistri}



\begin{quote}
Before the new one: say the Urdu for a customer, then the Urdu for business.
\end{quote}

\subsection*{You'll want to know: \ur{مستری}}
\textbf{\ur{مستری}} — \emph{mistrī} — \textquotedblleft{}a mechanic\textquotedblright{}.

\begin{cousinweb}
\textbf{\ur{مستری}} (\emph{mistrī}) is usually traced to Portuguese \emph{mestre}, \textquotedblleft{}master, master craftsman\textquotedblright{}.
\end{cousinweb}

\begin{grammarlens}[title={What you've built}]
One more word for this chapter.
\end{grammarlens}

\subsection*{Your turn}
\begin{itemize}
\raggedright
  \item \emph{Say it:} \emph{mistrī}
  \item \emph{Say it:} \emph{mistrī}, once more
  \item \emph{Say it:} say \emph{mistrī}
  \item \emph{Recall:} say the Urdu for a customer, then the Urdu for business, then say \emph{mistrī} again
\end{itemize}

\begin{tcolorbox}[breakable,colback=teal!4,colframe=teal!35!black,title={Before you move on}]
What does \ur{مستری} mean? (\textquotedblleft{}A mechanic\textquotedblright{}.) Say it once more.
\end{tcolorbox}
\section[\texorpdfstring{\ur{بیرا} (berā)}{بیرا (berā)}]{\ur{بیرا} (berā) — a waiter}

\label{lesson:UR-C237-bera}



\begin{quote}
Before the new one: say the Urdu for business, then the Urdu for a mechanic.
\end{quote}

\subsection*{You'll want to know: \ur{بیرا}}
\textbf{\ur{بیرا}} — \emph{berā} — \textquotedblleft{}a waiter\textquotedblright{}.

\begin{cousinweb}
\textbf{\ur{بیرا}} (\emph{berā}) is usually traced to English \emph{bearer}, a word the British in India used for a household servant.
\end{cousinweb}

\begin{grammarlens}[title={What you've built}]
One more word for this chapter.
\end{grammarlens}

\subsection*{Your turn}
\begin{itemize}
\raggedright
  \item \emph{Say it:} \emph{berā}
  \item \emph{Say it:} \emph{berā}, once more
  \item \emph{Say it:} say \emph{berā}
  \item \emph{Recall:} say the Urdu for business, then the Urdu for a mechanic, then say \emph{berā} again
\end{itemize}

\begin{tcolorbox}[breakable,colback=teal!4,colframe=teal!35!black,title={Before you move on}]
What does \ur{بیرا} mean? (\textquotedblleft{}A waiter\textquotedblright{}.) Say it once more.
\end{tcolorbox}
\section[\texorpdfstring{\ur{چوکیدار} (chaukīdār)}{چوکیدار (chaukīdār)}]{\ur{چوکیدار} (chaukīdār) — a watchman}

\label{lesson:UR-C237-chaukidar}



\begin{quote}
Before the new one: say the Urdu for a mechanic, then the Urdu for a waiter.
\end{quote}

\subsection*{You'll want to know: \ur{چوکیدار}}
\textbf{\ur{چوکیدار}} — \emph{chaukīdār} — \textquotedblleft{}a watchman\textquotedblright{}.

\begin{cousinweb}
\textbf{\ur{چوکیدار}} (\emph{chaukīdār}) is \textquotedblleft{}holder of the post\textquotedblright{}. \emph{chaukī}, \textquotedblleft{}a guard post\textquotedblright{}, is listed by CDIAL under Sanskrit \emph{catuṣka}, \textquotedblleft{}set of four; a square courtyard\textquotedblright{} (CDIAL 4629), and \emph{-dār}, \textquotedblleft{}holder\textquotedblright{}, is Persian.
\end{cousinweb}

\begin{grammarlens}[title={What you've built}]
One more word for this chapter.
\end{grammarlens}

\subsection*{Your turn}
\begin{itemize}
\raggedright
  \item \emph{Say it:} \emph{chaukīdār}
  \item \emph{Say it:} \emph{chaukīdār}, once more
  \item \emph{Say it:} say \emph{chaukīdār}
  \item \emph{Recall:} say the Urdu for a mechanic, then the Urdu for a waiter, then say \emph{chaukīdār} again
\end{itemize}

\begin{tcolorbox}[breakable,colback=teal!4,colframe=teal!35!black,title={Before you move on}]
What does \ur{چوکیدار} mean? (\textquotedblleft{}A watchman\textquotedblright{}.) Say it once more.
\end{tcolorbox}
\section[\texorpdfstring{\ur{سپاہی} (sipāhī)}{سپاہی (sipāhī)}]{\ur{سپاہی} (sipāhī) — a soldier; a constable}

\label{lesson:UR-C237-sipahi}



\begin{quote}
Before the new one: say the Urdu for a waiter, then the Urdu for a watchman.
\end{quote}

\subsection*{You'll want to know: \ur{سپاہی}}
\textbf{\ur{سپاہی}} — \emph{sipāhī} — \textquotedblleft{}a soldier; a constable\textquotedblright{}.

\begin{cousinweb}
\textbf{\ur{سپاہی}} (\emph{sipāhī}) is a loan, usually traced to Persian \emph{sipāhī}, \textquotedblleft{}soldier\textquotedblright{}. It is built on \emph{sipāh}, \textquotedblleft{}army\textquotedblright{}. English took the word as \emph{sepoy}.
\end{cousinweb}

\begin{grammarlens}[title={What you've built}]
One more word for this chapter.
\end{grammarlens}

\subsection*{Your turn}
\begin{itemize}
\raggedright
  \item \emph{Say it:} \emph{sipāhī}
  \item \emph{Say it:} \emph{sipāhī}, once more
  \item \emph{Say it:} say \emph{sipāhī}
  \item \emph{Recall:} say the Urdu for a waiter, then the Urdu for a watchman, then say \emph{sipāhī} again
\end{itemize}

\begin{tcolorbox}[breakable,colback=teal!4,colframe=teal!35!black,title={Before you move on}]
What does \ur{سپاہی} mean? (\textquotedblleft{}A soldier; a constable\textquotedblright{}.) Say it once more.
\end{tcolorbox}
\section[\texorpdfstring{\ur{فوج} (fauj)}{فوج (fauj)}]{\ur{فوج} (fauj) — an army}

\label{lesson:UR-C237-fauj}



\begin{quote}
Before the new one: say the Urdu for a watchman, then the Urdu for a soldier; a constable.
\end{quote}

\subsection*{You'll want to know: \ur{فوج}}
\textbf{\ur{فوج}} — \emph{fauj} — \textquotedblleft{}an army\textquotedblright{}.

\begin{cousinweb}
\textbf{\ur{فوج}} (\emph{fauj}) is a loan from Arabic \emph{fawj}, \textquotedblleft{}a troop, a company\textquotedblright{}, which reached Urdu through Persian. Its Arabic root is \emph{f-w-j} (Lane).
\end{cousinweb}

\begin{grammarlens}[title={What you've built}]
One more word for this chapter.
\end{grammarlens}

\subsection*{Your turn}
\begin{itemize}
\raggedright
  \item \emph{Say it:} \emph{fauj}
  \item \emph{Say it:} \emph{fauj}, once more
  \item \emph{Say it:} say \emph{fauj}
  \item \emph{Recall:} say the Urdu for a watchman, then the Urdu for a soldier; a constable, then say \emph{fauj} again
\end{itemize}

\begin{tcolorbox}[breakable,colback=teal!4,colframe=teal!35!black,title={Before you move on}]
What does \ur{فوج} mean? (\textquotedblleft{}An army\textquotedblright{}.) Say it once more.
\end{tcolorbox}
